% GPT-6 Astra Pro; /root/fixed_surface; 22 November 2026 UTC.
% Effort not exposed. Separate complete proof for review.
\section[HD-F1-S39: cohomology defects in codimension two]{HD-F1-S39: the normal obstruction survives adjacent cohomology maps supported in codimension at least two}
\label{sec:f1-norm-fourfold-codimension-two-cohomology}

Retain the actual normal integral norm image
$i:Y\hookrightarrow P$, its ideal $\mathcal I$,
$E=i_*\mathcal O_Y$, and the S29 geometric open.
Let $F\in D^b_{\rm coh}(P)$ be an actual complex,
write $C^a=H^a(F)$, and choose an integer $j$.
Its selected cohomology $M=C^j$ has set-theoretic
support $Y$ and the actual generic module
\[
 M_\xi\simeq\bigoplus_a
       (\mathcal O_{P,\xi}/\mathcal I_\xi^a)^{m_a},
 \qquad \ell_M=\sum_a m_a\binom{a+3}{4}>0.
 \tag{F1S39.1}
\]
There is a closed $Z\subset Y$ with
$\operatorname{codim}_Y Z\ge2$ such that every
ordinary global morphism in either of the groups
\[
 \operatorname{Hom}_P(C^{j-1},M),\qquad
 \operatorname{Hom}_P(M,C^{j+1})
 \tag{F1S39.2}
\]
has image supported in $Z$. Neither group is
required to vanish. Equivalently, their actual
restriction maps to $P\setminus Z$ are zero.
It also suffices that the image of each such
morphism have codimension at least two in $Y$:
the two Hom spaces are finite-dimensional, and
the union of the image supports of finite bases
gives one such $Z$.

We prove the actual characteristic-image surjection
\[
\begin{split}
 \mathcal T_F:\operatorname{ev}_F
   (\mathfrak D_P\oplus H^2(P,\mathcal O_P))
       &\longrightarrow W_E,\\
 \operatorname{ev}_F(\kappa+\eta)
       &\longmapsto-\ell_M\operatorname{ob}_E(\kappa),
 \qquad W_E=\operatorname{ob}_E(\mathfrak D_P).
\end{split}
 \tag{F1S39.3}
\]
At the complete S14 character this gives an
invariant full degree-two trace kernel of dimension
at least $16-r_C\ge4$ for every specified coherent
finite point lift. Paragraphs 5--7 construct an
actual gluable, ordinarily indecomposable complex
with the complete character, nonzero group
(F1S39.2), and no $Ri_*D^b_{\rm coh}(Y)$
presentation.

\paragraph{1. The actual lifted cohomology is flat on the chosen open.}
Put $V=P\setminus Z$ and $U=Y\setminus Z$.
Suppose $\operatorname{ev}_F(\kappa+\eta)=0$.
The precise Toda/BF comparison retained in
S33--S38 supplies an actual scalar-twisted
perfect lift $\widetilde F$ on the ambient
first-order deformation $P_\kappa$.
For its central immersion $s:P\hookrightarrow
P_\kappa$, the actual epsilon triangle is
\[
 s_*F\xrightarrow{u}\widetilde F
       \xrightarrow{v}s_*F\longrightarrow s_*F[1],
 \qquad vu=0,\qquad uv=\epsilon\operatorname{id}.
 \tag{F1S39.4}
\]
On a local bounded finite free representative it
is the literal sequence of complexes
$0\to\epsilon L^\bullet\to L^\bullet
\to L^\bullet/\epsilon L^\bullet\to0$,
with the first and last terms identified by
multiplication by epsilon. Thus every differential
and every coupling of the given lift is retained.

The connecting cohomology maps are ordinary
global $\mathcal O_P$-linear maps
$b_a:C^a\to C^{a+1}$. Indeed they are the
connecting maps of (F1S39.4); its first and last
terms are annihilated by epsilon. On a triple
overlap the scalar $1-\epsilon\eta_{ikl}$
therefore acts as the identity on both terms,
so their cohomology maps descend as ordinary
maps on $P$. Both $b_{j-1}$ and $b_j$ vanish
after restriction to $V$ by (F1S39.2).
The exact cohomology sequence consequently gives
\[
 0\longrightarrow s_*(M|_V)\xrightarrow{u_j}
 H^j(\widetilde F)|_V\xrightarrow{v_j}
 s_*(M|_V)\longrightarrow0,
 \qquad u_jv_j=\epsilon.
 \tag{F1S39.5}
\]
Write $\widetilde M=H^j(\widetilde F)|_V$.
Surjectivity of $v_j$ gives
$\operatorname{im}\epsilon=\operatorname{im}u_j$.
If $\epsilon x=0$, then $u_jv_jx=0$;
injectivity of $u_j$ gives $v_jx=0$ and hence
$x\in\operatorname{im}u_j$. Therefore
\[
 \ker(\epsilon:\widetilde M\to\widetilde M)
       =\epsilon\widetilde M,
 \qquad \widetilde M/\epsilon\widetilde M=M|_V.
 \tag{F1S39.6}
\]
These identities prove epsilon-flatness directly.
At a stalk, lift a complex-vector-space basis of
the quotient. It generates over the dual numbers
by $Q=\epsilon Q\Rightarrow Q=0$.
A relation first reduces to zero coefficients;
dividing its remaining epsilon coefficients and
using (F1S39.6) again makes them zero.
The resulting module is free over the dual
numbers, possibly of infinite vector-space rank.
The actual cohomology overlap maps retain the
same scalar triple products as the perfect lift.
No flatness is asserted on $Z$ or in any other
cohomological degree.

\paragraph{2. Local normal maps and every ideal-generator extension.}
Choose a finite affine cover $V_i$ of $P$ that
trivializes the smooth ambient deformation.
Restrict the flat sheaf (F1S39.5) to
$V_i\setminus Z$. The complete local
flat-sheaf construction in S37 paragraphs 1--3
applies to this actual sheaf: its central
generic module is (F1S39.1), all off-diagonal
first-order actions are retained, and its
epsilon-flat coefficient-ring resolution gives
regular values at every smooth support point
of $U$. The normal-domain extension used there
also applies at its remaining singular points.
It gives actual maps
\[
 u_i^\circ:(\mathcal I/\mathcal I^2)|_{U_i}
                 \longrightarrow\mathcal O_{U_i},
 \qquad U_i=(V_i\cap Y)\setminus Z.
 \tag{F1S39.7}
\]
This invocation uses only the local construction
and its given flat module. It makes no assumption
that the whole $\widetilde F$ or its other
cohomology sheaves are flat, and no choice of a
generic coefficient field for a coupled complex
enters it.

Here is the exact extension through the omitted
set. Put $R_i=\Gamma(V_i\cap Y,\mathcal O_Y)$,
a normal domain when this affine intersection
is nonempty. Choose actual generators $f_{is}$
of the ambient ideal on $V_i$. Each value
$u_i^\circ(f_{is})$ is a rational function on
$Y$ regular on $U_i$. Every height-one point
of $\operatorname{Spec}R_i$ is in $U_i$.
The actual normal-domain intersection gives
\[
 \Gamma(U_i,\mathcal O_{U_i})
       \longrightarrow\operatorname{Frac}(R_i),
 \qquad
 R_i=\bigcap_{\operatorname{ht}\mathfrak p=1}
                 (R_i)_{\mathfrak p}.
 \tag{F1S39.8}
\]
Thus each generator value has a unique extension
$b_{is}\in R_i$. If
$\sum_s a_s f_{is}\in I^2$, its evaluation
$\sum_s\overline a_s b_{is}$ vanishes on
$U_i$ and hence in the domain $R_i$.
All original relations, including those killing
$I^2$, therefore remain valid. We obtain the
actual $\mathcal O_Y$-linear morphism
\[
 u_i:(\mathcal I/\mathcal I^2)|_{V_i\cap Y}
                   \longrightarrow\mathcal O_{V_i\cap Y}.
 \tag{F1S39.9}
\]
Uniqueness on generators shows independence of
the presentation and compatibility with every
further affine restriction. This is the
literal Hartogs restriction-and-extension map;
no extension of a generic nonreduced matrix
trace is asserted.

\paragraph{3. Ambient transitions, scalar twists and flat ideals.}
Orient the local trivialization transition
from chart $i$ to chart $k$ as
$\phi_{ik}=1+\epsilon\kappa_{ik}$.
On $U_i\cap U_k$, the exact S37 overlap
calculation for the given flat sheaf gives
\[
              u_i^\circ-u_k^\circ
                     =\partial_1(\kappa_{ik}).
 \tag{F1S39.10}
\]
It is obtained from conjugation of the actual
diagonal ideal-action maps and the identity
$\rho_a\partial_a=\partial_1$; their
length-weighted sum has total coefficient one.
The scalar triple products from $\eta$
act trivially by conjugation. After evaluating
on any local ideal generator, both sides of
(F1S39.10) are regular on the entire overlap
in $Y$ by (F1S39.9). Their difference vanishes
on its dense complement of $Z$, so it is zero.
Thus the original ambient transition and
scalar-twist identities descend through exactly
the same Hartogs maps, with all relations retained.

Define actual local ideals in the trivial
ambient deformations by
\[
 I_{\epsilon,i}=\{f+\epsilon h:
          f\in I_i,\quad h\bmod I_i=-u_i(f)\}.
 \tag{F1S39.11}
\]
Conormal linearity makes them ideals.
Their quotient reduces to $R/I_i$.
If epsilon kills the class of $r+\epsilon s$,
then $r\in I_i$. Subtracting the displayed
lift of $r$ leaves an epsilon class, proving
$\ker\epsilon=\operatorname{im}\epsilon$.
They are therefore epsilon-flat by the direct
criterion proved after (F1S39.6).
The transition $\phi_{ik}$ takes the epsilon
coefficient $-u_i(f)$ to
$-u_i(f)+\kappa_{ik}(f)=-u_k(f)$.
It carries $I_{\epsilon,i}$ to $I_{\epsilon,k}$.
These ideals glue in the actual $P_\kappa$
and deform the original reduced embedded $Y$.
The S24 obstruction comparison consequently gives
\[
 \operatorname{ev}_F(\kappa+\eta)=0
              \quad\Longrightarrow\quad
                        \operatorname{ob}_E(\kappa)=0.
 \tag{F1S39.12}
\]
Kernel inclusion proves that (F1S39.3) is
well defined and linear, by subtraction of
representatives. Its surjectivity follows from
$\ell_M>0$ and the definition of $W_E$.
The minus sign is the fixed S31 first-jet/HKR
dictionary. The coefficient is the selected
sheaf's length, not an asserted Euler length
of the entire complex.

\paragraph{4. The complete character and the specified point action.}
Retain the full equality on the actual complex
\[
 e^B\operatorname{ch}(F)=\lambda t+Q(\beta),
 \quad B\in H^{1,1}(P,\mathbf Q),\quad
 0\ne\lambda t\in W_K(P),\quad Q\in\mathbf Q[x].
 \tag{F1S39.13}
\]
The full BF exponential derivative, through
the exact S31--S38 source maps, gives zero
full trace on
\[
 \operatorname{ev}_F
   \{\kappa+\kappa\mathbin{\lrcorner}B:
                            \kappa\in\mathfrak D_P\}.
 \tag{F1S39.14}
\]
Its image under (F1S39.3) is $W_E$.
Every specified coherent finite point lift
fixes (F1S39.14) by the original
Hochschild-action naturality: translations and
Picard-zero tensoring fix this source.
The actual R05 trace comparison therefore gives
\[
             \dim\ker\sigma_F^2\ge16-r_C\ge4
 \tag{F1S39.15}
\]
in the invariant full-trace problem. All
higher character terms and all original
couplings remain in (F1S39.13).

\paragraph{5. Actual point ideals and a nonzero adjacent Hom group.}
Choose distinct $p,q,r\in Y_{\rm reg}$.
Put $J_s=I_{s/Y}$, $E_s=i_*J_s$, and
retain the exact maps
\[
 0\longrightarrow E_s\xrightarrow{a_s}E
                   \longrightarrow k_s\longrightarrow0.
 \tag{F1S39.16}
\]
Set $A=E_p\oplus E_r$ and $T=E_q\oplus k_p$.
The exact rational-multiplication and
height-one extension argument of S38 gives
$\operatorname{End}(E_s)=\mathbf C$ and
$\operatorname{Hom}(E_s,E_t)=0$ for $s\ne t$.
There are no maps from a point sheaf to these
torsion-free sheaves on integral $Y$.
Consequently
\[
 \operatorname{Hom}_P(T,A)=0,\qquad
 \operatorname{Hom}_P(A,T)
     =\operatorname{Hom}_P(E_p\oplus E_r,k_p)
     \simeq\mathbf C^4\oplus\mathbf C.
 \tag{F1S39.17}
\]
The four-dimensional term is the dual of
$\mathfrak m_{Y,p}/\mathfrak m_{Y,p}^2$;
the other point is disjoint from $r$.
Thus every forward map lands in the actual
subsheaf $k_p\subset T$, and its image is
supported in $Z=\{p\}$, whereas this Hom
group is nonzero.

For $s=p,r$, S38 supplies actual isomorphisms
and the normal comparison
\[
\begin{split}
 a_q^*:\operatorname{Ext}^2_P(E,E_s)
       &\xrightarrow{\sim}\operatorname{Ext}^2_P(E_q,E_s),\\
 \Psi_{q,s}&=(a_s)_*(a_q^*)^{-1},
 \qquad \Psi_{q,s}(\delta_s)=w_s\ne0,
       \quad w_s\in W_E.
\end{split}
 \tag{F1S39.18}
\]
Choose two nonzero such classes $w_p,w_r$
and their actual lifts $\delta_p,\delta_r$.
Their existence uses the local-to-global edge
zero of $W_E$, the actual target point-ideal
sequence and the source isomorphism above,
exactly as in S38; it is not an assumption
that a whole-character class lifts to an object.

At the smooth point $p$, the actual ambient
Koszul calculation gives
$\operatorname{Ext}^a_P(k_p,E)=0$ for $a<4$.
The connecting map of (F1S39.16) is therefore
an isomorphism
\[
 \operatorname{Ext}^1_P(k_p,k_p)
       \xrightarrow{\sim}\operatorname{Ext}^2_P(k_p,E_p).
 \tag{F1S39.19}
\]
Its source has dimension eight by the
residue-field Koszul complex in the smooth
eight-dimensional ambient local ring.
Choose any nonzero class $\tau$ in its image.
Its pushforward by $a_p$ is zero, by the same
exact sequence. Since $p\ne r$,
$\operatorname{Ext}^2_P(k_p,E_r)=0$.
Thus the actual matrix
\[
 \delta:T\longrightarrow A[2],\qquad
 \delta=\begin{pmatrix}\delta_p&\tau\\
                         \delta_r&0\end{pmatrix}
 \tag{F1S39.20}
\]
is defined, with every displayed nonzero
entry supplied by an actual Ext group.
Take its actual cone in the triangle
$A\xrightarrow{i}F\xrightarrow{p_F}T[-1]
\xrightarrow{\delta[-1]}A[1]$.
It has $H^0(F)=A$, $H^1(F)=T$, and no
other cohomology. The selected $M=A$
has generic module $K(Y)^2$, so $\ell_M=2$.
Equations (F1S39.17) verify (F1S39.2)
for $j=0$ with this nonempty point support.

\paragraph{6. Gluability and the full nilpotent endomorphism ideal.}
The two exact Hom sequences for the displayed
triangle give
$\operatorname{Ext}^{-1}_P(F,F)
\simeq\operatorname{Hom}_P(T,A)$ by
$h\mapsto i[-1]h[-1]p_F$.
All preceding negative Ext groups of the
coherent sheaves vanish. For degrees at most
$-2$ the usual cohomological bounds give zero
directly. Hence this actual $F$ is gluable.

The ordinary endomorphism rings of its
cohomology have matrices
\[
 \operatorname{End}(A)=\mathbf C^2,
 \qquad
 \operatorname{End}(T)=
  \left\{\begin{pmatrix}c&0\\v&t\end{pmatrix}:
        c,t\in\mathbf C,\quad
        v\in\operatorname{Hom}(E_q,k_p)\right\}.
 \tag{F1S39.21}
\]
Here $\operatorname{Hom}(E_q,k_p)=\mathbf C$.
For an actual endomorphism of $F$, write its
two scalars on $A$ as $\alpha_p,\alpha_r$.
Naturality of its truncation triangle gives
\[
 \alpha_p\delta_p=c\delta_p+\tau v,
 \qquad \alpha_r\delta_r=c\delta_r,
 \qquad \alpha_p\tau=t\tau.
 \tag{F1S39.22}
\]
The middle equality and $\delta_r\ne0$
give $\alpha_r=c$. The exact comparison
$\Psi_{q,p}$ kills $\tau v$: its target
pushforward by $a_p:E_p\to E$ is zero.
The source pullback $a_q^*$ is also an
isomorphism with target $E$, by the same
vanishings $\operatorname{Ext}^2_P(k_q,E)
=\operatorname{Ext}^3_P(k_q,E)=0$.
The resulting commutative square with target
pushforward therefore proves the asserted
vanishing under $\Psi_{q,p}$.
Applying (F1S39.18) to the first equality
gives $\alpha_p=c$.
Finally $\tau\ne0$ gives $t=c$.
Thus the image of
\[
 \operatorname{End}(F)\longrightarrow
                 \operatorname{End}(A)\oplus\operatorname{End}(T)
 \tag{F1S39.23}
\]
lies in the scalars plus the strictly lower
off-diagonal subspace in (F1S39.21).
That latter subspace has square zero; the
additional compatibility $\tau v=0$ is
retained, not discarded.

The kernel $K_0$ of (F1S39.23) also has
square zero, by exact maps as follows.
The cohomological bounds give
$\operatorname{Hom}(A,F)=\operatorname{End}(A)$
through $i$, and
$\operatorname{Hom}(F,T[-1])=\operatorname{End}(T)$
through $p_F$. Hence an element $e\in K_0$
satisfies $ei=0$ and $p_Fe=0$.
The first equality factors it as $e=h p_F$.
Precomposition by $p_F$ on
$\operatorname{End}(T[-1])$ is injective,
so $p_Fh=0$. The second exact Hom sequence
then gives $h=i v_1$ with
$v_1\in\operatorname{Ext}^1_P(T,A)$.
Every element of $K_0$ is thus $i v_1p_F$,
and two such elements multiply to zero
because $p_Fi=0$.

Let $J_0$ be the kernel of the common-scalar
map $\operatorname{End}(F)\to\mathbf C$
just obtained. Its quotient by $K_0$ is
contained in the square-zero off-diagonal
space, so $J_0^2\subset K_0$ and
$J_0^4=0$. An idempotent has scalar zero
or one. In the first case it is in this
nilpotent ideal and hence zero; in the
second case apply the same argument to
its complement. There are no nontrivial
ordinary idempotents, proving ordinary
indecomposability without an assumption
that either endomorphism projection is onto.

\paragraph{7. Actual support comparison and the whole character.}
This complex has no $Ri_*G$ presentation
with $G\in D^b_{\rm coh}(Y)$. Indeed,
exact closed pushforward commutes with
cohomology and truncation. Its ordinary
Hom map is fully faithful, so an alleged
presentation supplies the actual cohomology
identifications on $Y$ and places the
coupling (F1S39.20) in
\[
 \operatorname{im}\left[
   \operatorname{Ext}^2_Y(J_q\oplus k_p,J_p\oplus J_r)
      \longrightarrow\operatorname{Ext}^2_P(T,A)\right].
 \tag{F1S39.24}
\]
This image is stable under the actual
ordinary cohomology identifications.
Projection to the $J_q\to J_p$ entry
would put $\delta_p$ in the corresponding
support Ext image. S38's exact source
point-ideal isomorphism and target pushforward
send that image under $\Psi_{q,p}$ to the
scalar $H^2(Y,\mathcal O_Y)$ image in
$\operatorname{Ext}^2_P(E,E)$.
The actual scalar retraction is the identity
there and zero on $W_E$, whereas
$\Psi_{q,p}(\delta_p)=w_p\ne0$.
This contradicts (F1S39.24) through the
displayed maps.

The actual K-theory equality is
$[F]=[E]-2[k_p]-[k_r]+[k_q]$.
Apply the cycle-class map to its Chern
character. On the connected abelian
eightfold all point cohomology classes
equal $\beta^8/\int_P\beta^8$, so
\[
 \operatorname{ch}(F)=\operatorname{ch}(E)
                  -2\frac{\beta^8}{\int_P\beta^8}.
 \tag{F1S39.25}
\]
No rational equivalence between different
points is asserted in the Chow group.
Equation (F1S39.25) is the full S14
character with $B=0$, the same nonzero
determinant component, and the displayed
top polynomial correction. The trivial
point group gives an actual coherent lift;
every other specified lift that exists
is covered by (F1S39.15).
The nonzero group (F1S39.17) verifies that
this actual complex is outside S38's
adjacent-Hom-vanishing hypothesis.

\paragraph{Source scope and exact remaining boundary.}
The only new extension operation is the
displayed Hartogs map (F1S39.8) on the
normal reduced support, with every generator
relation and ambient transition checked above.
The complete primary Stacks 0AVB proof was
freshly read; the complete 0149 proof was
also freshly read for every exact Hom
factorization used here. The actual epsilon
triangle, Toda scalar-twisted lift, corrected
jet/HKR sign, flat-sheaf normal map and
point-ideal comparisons retain the completely
reviewed S33--S38 scopes. The point Ext
groups used in (F1S39.17)--(F1S39.19) are
computed by the literal regular-local Koszul
complexes, with faithful completion and
point-supported local-to-global comparison
as in S38. No unread assertion about arbitrary
derived coefficient-field invariance is used.

The proof does not require the adjacent
Bocksteins to vanish on $Z$, nor assert
cohomology-flatness there. Adjacent maps
whose image survives at a height-one point
of $Y$ remain outside this result unless
another actual argument annihilates the
corresponding Bocksteins. Arbitrary generic
thickening modules also retain the exact
S37 boundary. No assertion about either
remaining gate is inferred from Hartogs
extension.
